\documentclass[12pt,a4paper]{article}

% ── Packages ──────────────────────────────────────────────
\usepackage[utf8]{inputenc}
\usepackage[T1]{fontenc}
\usepackage{lmodern}
\usepackage{amsmath,amssymb,amsthm}
\usepackage{mathtools}
\usepackage{geometry}
\geometry{margin=1in}
\usepackage{graphicx}
\graphicspath{{../figures/}}
\usepackage{booktabs}
\usepackage{hyperref}
\hypersetup{
    colorlinks=true,
    linkcolor=blue!60!black,
    citecolor=blue!60!black,
    urlcolor=blue!60!black
}
\usepackage{caption}
\usepackage{float}
\usepackage{enumitem}
\usepackage{xcolor}

% ── Custom commands ───────────────────────────────────────
\newcommand{\Hfour}{H_4}
\newcommand{\Hthree}{H_3}
\newcommand{\Htwo}{H_2}
\newcommand{\Eeight}{E_8}
\newcommand{\vphi}{\varphi}
\newcommand{\Lmat}{L_{H_4}}
\newcommand{\R}{\mathbb{R}}
\newcommand{\Q}{\mathbb{Q}}
\newcommand{\Z}{\mathbb{Z}}

\newtheorem{invariant}{Invariant}
\newtheorem{proposition}{Proposition}

% ── Title ─────────────────────────────────────────────────
\title{%
  \textbf{Algebraic Invariants of the $\Hfour$ Laplacian Spectrum}\\[6pt]
  \large Part~II of a Series on Symmetry-Constrained Resonant Structure%
}
\author{%
  Lee Smart\\[4pt]
  \textit{Institute of Vibrational Field Dynamics}\\[2pt]
  \texttt{contact@vibrationalfielddynamics.org}\\[2pt]
  \texttt{@vfd\_org}%
}
\date{\today}

% ══════════════════════════════════════════════════════════
\begin{document}
\maketitle
\thispagestyle{empty}

% ── Abstract ──────────────────────────────────────────────
\begin{abstract}
The graph Laplacian of the 600-cell --- the terminal regular polytope
of the non-crystallographic Coxeter family~$H_n$ --- possesses a
spectrum of exceptional algebraic regularity.  Across 120~degrees of
freedom, only nine distinct eigenvalues appear.  The associated
multiplicities include the perfect-square sequence
$(1,4,9,16,25,36)$, the irrational eigenvalues occur in Galois
conjugate pairs anchored at the graph degree~$k=12$, and the total
$\vphi$-content of the spectrum cancels exactly.  This paper
identifies and formalizes seven structural invariants governing this
spectrum, relates them to the harmonic structure of the 3-sphere~$S^3$,
and discusses their algebraic origin in terms of the distance-regular
geometry and symmetry algebra of the 600-cell graph.  These invariants
are shown to be consistent with the coherent nonlinear dynamics
reported in Part~I of this series.
\end{abstract}

\vspace{1em}
\noindent\textbf{Keywords:}
Coxeter groups, $\Hfour$ symmetry, 600-cell, graph Laplacian,
spectral invariants, algebraic graph theory, golden ratio,
hyperspherical harmonics.

\newpage
\tableofcontents
\newpage


% ══════════════════════════════════════════════════════════
\section{Introduction}\label{sec:intro}
% ══════════════════════════════════════════════════════════

In Part~I of this series~\cite{paper1}, the 600-cell was introduced as a
geometric substrate for a nonlinear field model constrained by the
symmetry of the Coxeter group~$\Hfour$.  The Laplacian spectrum of the
600-cell graph was computed numerically and found to contain only nine
distinct eigenvalues across 120~degrees of freedom.  Numerical
simulations demonstrated that the spectral compression produces
persistent localisation, coherent shell reflection, and recurrent
breathing dynamics that are absent in degree-matched control graphs.

The present paper studies the internal organization of this
nine-sector Laplacian spectrum.  Rather than treating the eigenvalues
as numerical data, we analyze their algebraic structure and identify a
set of seven structural invariants that govern the spectrum.

These invariants are not assumed; they are extracted from the computed
spectrum and verified exactly.  They concern:
\begin{enumerate}[nosep]
\item the cancellation of golden-ratio content across eigenvalues,
\item the anchoring of irrational eigenvalues at the graph degree,
\item the pairing structure of Galois conjugate eigenvalues,
\item the folding of multiplicities from a perfect-square sequence,
\item the partition of modes into three $\vphi$-zones,
\item the decomposition of the spectral backbone,
\item the structure of the discrete departure from spherical harmonics.
\end{enumerate}

These invariants are mutually consistent and arise simultaneously from
the same spectral data; none is imposed independently.

Together these invariants suggest that the 600-cell Laplacian
spectrum is far more algebraically constrained than a generic
graph spectrum, and they are consistent with the dynamical phenomena
reported in Part~I.

For convenience, these invariants will sometimes be referred to
informally as the ``seven laws'' of the $\Hfour$ spectrum, without
implying independence or axiomatic status.

The focus of the present work is strictly algebraic.  No physical
interpretation is assumed.  The goal is to identify structural
constraints of the $\Hfour$ Laplacian spectrum and to determine whether
these constraints arise from known algebraic or geometric principles.


% ══════════════════════════════════════════════════════════
\section{Spectral Data of the 600-Cell Graph}\label{sec:data}
% ══════════════════════════════════════════════════════════

The 600-cell is a regular convex 4-polytope with 120~vertices,
720~edges, 1200~triangular faces, and 600~tetrahedral cells.  Its
graph is 12-regular and vertex-transitive under the Coxeter
group~$\Hfour$ of order~14\,400.

The graph Laplacian is defined as $\Lmat = 12I - A$, where $A$ is
the $120 \times 120$ adjacency matrix.  The eigenvalues of~$\Lmat$
are listed in Table~\ref{tab:spectrum}.

\begin{table}[H]
\centering
\caption{Eigenvalues and multiplicities of the 600-cell graph
Laplacian $\Lmat = 12I - A$.}
\label{tab:spectrum}
\begin{tabular}{@{}cllrl@{}}
\toprule
Sector & Exact form & $\vphi$-basis ($a+b\vphi$) & Mult.\ $d_k$ & Note \\
\midrule
1 & $0$ & $0 + 0\vphi$ & 1 & $= 1^2$ \\
2 & $9 - 3\sqrt{5}$ & $12 - 6\vphi$ & 4 & $= 2^2$ \\
3 & $10 - 2\sqrt{5}$ & $12 - 4\vphi$ & 9 & $= 3^2$ \\
4 & $9$ & $9 + 0\vphi$ & 16 & $= 4^2$ \\
5 & $12$ & $12 + 0\vphi$ & 25 & $= 5^2$ \\
6 & $14$ & $14 + 0\vphi$ & 36 & $= 6^2$ \\
7 & $10 + 2\sqrt{5}$ & $8 + 4\vphi$ & 9 & \\
8 & $15$ & $15 + 0\vphi$ & 16 & \\
9 & $9 + 3\sqrt{5}$ & $6 + 6\vphi$ & 4 & \\
\bottomrule
\addlinespace
\multicolumn{5}{@{}l}{\footnotesize
  $\sum d_k = 120$.  Conversion: $\sqrt{5} = 2\vphi - 1$.
  All eigenvalues lie in $\Q(\sqrt{5}\,)$.}
\end{tabular}
\end{table}

The spectrum is computed exactly and all algebraic expressions are
verified symbolically within~$\Q(\sqrt{5}\,)$.  All multiplicities
and algebraic identities were cross-checked against numerical
eigendecomposition to machine precision.

Every eigenvalue can be written in the form $\mu = a + b\vphi$ where
$a \in \Z$ and $b \in \{-6,-4,0,4,6\}$.  The four irrational
eigenvalues form two Galois conjugate pairs under the automorphism
$\sqrt{5} \mapsto -\sqrt{5}$:
\begin{align}
  \text{Pair 1:} &\quad (9 - 3\sqrt{5},\; 9 + 3\sqrt{5})
    \qquad \text{sum} = 18,\; \text{product} = 36, \label{eq:pair1} \\
  \text{Pair 2:} &\quad (10 - 2\sqrt{5},\; 10 + 2\sqrt{5})
    \qquad \text{sum} = 20,\; \text{product} = 80. \label{eq:pair2}
\end{align}


% ══════════════════════════════════════════════════════════
\section{The Seven Structural Invariants}\label{sec:invariants}
% ══════════════════════════════════════════════════════════

We now state the seven structural invariants of the spectrum.  Each
invariant is presented with an observation, exact statement, minimal
statement, interpretation, derivation status, type classification, and
implication.

% ── Invariant 1 ──────────────────────────────────────────
\begin{invariant}[$\vphi$-Cancellation]\label{inv:phi}

\emph{Observation.}
The irrational components of the spectrum appear symmetrically:
$9 \pm 3\sqrt{5}$ and $10 \pm 2\sqrt{5}$, and cancel when summed.

\emph{Exact statement.}
Decomposing each eigenvalue as $\mu_k = \mu_k^{(\Q)} + \mu_k^{(\sqrt{5})}$
into rational and irrational parts,
\begin{equation}\label{eq:phi_cancel}
  \sum_{k=1}^{9} \mu_k^{(\sqrt{5})} = 0
  \qquad\text{and}\qquad
  \sum_{k=1}^{9} d_k \, \mu_k^{(\sqrt{5})} = 0.
\end{equation}

\emph{Verification.}
The irrational parts are $\{-3\sqrt{5},\; -2\sqrt{5},\; 0,\; 0,\; 0,\;
0,\; 0,\; +2\sqrt{5},\; +3\sqrt{5}\}$, summing to zero.
Multiplicity-weighted: $(-3\sqrt{5})(4) + (-2\sqrt{5})(9) +
(2\sqrt{5})(9) + (3\sqrt{5})(4) = 0$.

\emph{Minimal statement.}
$\sum b_k = 0$ and $\sum b_k d_k = 0$.

\emph{Interpretation.}
The spectrum is invariant under the Galois conjugation
$\sqrt{5} \mapsto -\sqrt{5}$.  This is consistent with the Laplacian
spectrum lying in the number field~$\Q(\sqrt{5}\,)$ and being closed
under its Galois group.  The cancellation is a field symmetry
constraint, not a numerical coincidence.

\emph{Status:} Exact (algebraic). \quad
\textbf{Type:} Galois symmetry constraint.

\emph{Implication.}
The spectrum is closed under algebraic conjugation, restricting
the set of possible eigenvalue structures.
\end{invariant}

% ── Invariant 2 ──────────────────────────────────────────
\begin{invariant}[Degree Anchoring]\label{inv:degree}

\emph{Observation.}
The graph degree $k = 12$ appears as an eigenvalue and serves as the
central reference point of the spectrum.

\emph{Exact statement.}
The Laplacian is defined by $\Lmat = kI - A$ with $k = 12$.  The
value $\mu = 12$ is one of the nine distinct eigenvalues.  Furthermore,
for every eigenvalue with negative $\vphi$-coefficient ($b_k < 0$),
the rational part satisfies $a_k = k = 12$.

\emph{Verification.}
Sector~2: $\mu = 12 - 6\vphi$, $a = 12$.
Sector~3: $\mu = 12 - 4\vphi$, $a = 12$.

\emph{Minimal statement.}
$b_k < 0 \implies a_k = 12$.

\emph{Interpretation.}
The degree acts as a spectral anchor.  All irrational eigenvalues
are organised relative to this fixed combinatorial constraint.  The
$k = 12$ regularity of the graph determines the reference frame of
the spectrum.

\emph{Status:} Exact (follows from Laplacian definition and spectral
data). \quad \textbf{Type:} Algebraic (combinatorial anchor).

\emph{Implication.}
The spectrum is organised relative to the fixed combinatorial
constraint $k = 12$.
\end{invariant}

% ── Invariant 3 ──────────────────────────────────────────
\begin{invariant}[Conjugate Pairing]\label{inv:conj}

\emph{Observation.}
All irrational eigenvalues appear in Galois conjugate pairs with equal
multiplicities.

\emph{Exact statement.}
If $\lambda = a + b\sqrt{5}$ appears in the spectrum with
multiplicity~$d$, then $\lambda' = a - b\sqrt{5}$ also appears with
multiplicity~$d$.

\emph{Verification.}
Pair~1: $9 - 3\sqrt{5}$ and $9 + 3\sqrt{5}$, both multiplicity~4.
Pair~2: $10 - 2\sqrt{5}$ and $10 + 2\sqrt{5}$, both multiplicity~9.

\emph{Minimal statement.}
$\mu \in \mathrm{Spec}(\Lmat)$ irrational $\implies$
$\bar\mu \in \mathrm{Spec}(\Lmat)$ with $d(\mu) = d(\bar\mu)$.

\emph{Interpretation.}
The spectrum respects the Galois symmetry of
$\Q(\sqrt{5}\,)/\Q$.  This is consistent with the underlying algebra
generating the spectrum being defined over~$\Q(\sqrt{5}\,)$ and
with the $\vphi$-geometry of the non-crystallographic Coxeter family.

\emph{Status:} Exact (spectral property). \quad
\textbf{Type:} Galois symmetry constraint.

\emph{Implication.}
Eigenvalues are constrained by number field symmetry, reducing the
set of possible spectral forms.
\end{invariant}

% ── Invariant 4 ──────────────────────────────────────────
\begin{invariant}[Square Multiplicity Folding]\label{inv:fold}

\emph{Observation.}
The first six multiplicities follow a perfect-square sequence; the
upper three fold back to earlier values.

\emph{Exact statement.}
\begin{equation}\label{eq:fold}
  d_\ell = (\ell+1)^2, \qquad \ell = 0,1,\ldots,5,
\end{equation}
giving $d_1,\ldots,d_6 = 1, 4, 9, 16, 25, 36$.  The remaining
multiplicities are
\begin{equation}\label{eq:upper}
  d_7 = 9 = 3^2, \quad d_8 = 16 = 4^2, \quad d_9 = 4 = 2^2.
\end{equation}

\emph{Verification.}
$d_7 = 9 = d_3$, $d_8 = 16 = d_4$, $d_9 = 4 = d_2$.

\emph{Minimal statement.}
$d_k = (k)^2$ for $k = 1,\ldots,6$; \;
$d_7 = d_3$, $d_8 = d_4$, $d_9 = d_2$.

\emph{Interpretation.}
The $(\ell+1)^2$ pattern matches the degeneracies of hyperspherical
harmonics on~$S^3$.  This is consistent with the interpretation that
the 600-cell graph encodes a finite, symmetry-reduced discrete
analogue of~$S^3$ harmonic structure.  The upper multiplicities
represent a fold-back where the finite graph cannot support higher
harmonics.

\emph{Status:} Empirical (strong structural match; derivation from
$\Hfour$ representation theory is an open problem). \quad
\textbf{Type:} Combinatorial (multiplicity pattern).

\emph{Implication.}
The 600-cell graph appears to encode a symmetry-reduced discrete
analogue of~$S^3$ harmonic structure.  No comparable square-folding
pattern is observed in generic regular graphs of similar size and
degree.
\end{invariant}

% ── Invariant 5 ──────────────────────────────────────────
\begin{invariant}[Three-Zone Partition]\label{inv:zones}

\emph{Observation.}
The spectrum splits into three zones by $\vphi$-content, with
symmetric outer zones.

\emph{Exact statement.}
Define the spectral partition
\begin{equation}\label{eq:zones}
  \mathcal{S} =
  \mathcal{S}_{-} \cup \mathcal{S}_{0} \cup \mathcal{S}_{+}
\end{equation}
where $\mathcal{S}_{-}$ contains eigenvalues with $b < 0$,
$\mathcal{S}_{0}$ contains those with $b = 0$, and
$\mathcal{S}_{+}$ contains those with $b > 0$.  Then:
\begin{align}
  |\mathcal{S}_{-}| &= 13 \text{ modes}
    \quad (2 \text{ eigenvalues}), \\
  |\mathcal{S}_{0}| &= 94 \text{ modes}
    \quad (5 \text{ eigenvalues}), \\
  |\mathcal{S}_{+}| &= 13 \text{ modes}
    \quad (2 \text{ eigenvalues}).
\end{align}

\emph{Verification.}
$13 + 94 + 13 = 120$.  $|\mathcal{S}_{-}| = |\mathcal{S}_{+}| = 13$.

\emph{Minimal statement.}
$|\mathcal{S}_{-}| = |\mathcal{S}_{+}| = 13$; \;
$|\mathcal{S}_{0}| = 94$; \; $13 + 94 + 13 = 120$.

\emph{Interpretation.}
The spectrum is not uniform but organised into symmetry-related
sectors.  The outer zones behave symmetrically, carrying the
$\vphi$-geometry.  The middle zone carries the bulk rational
structure.  This stratification is consistent with the Galois
symmetry of~$\Q(\sqrt{5}\,)/\Q$.

\emph{Status:} Exact (direct count; formal derivation from
representation theory is an open problem). \quad
\textbf{Type:} Algebraic (Galois partition).

\emph{Implication.}
The spectrum is organised into symmetry-related sectors rather
than distributed uniformly.
\end{invariant}

% ── Invariant 6 ──────────────────────────────────────────
\begin{invariant}[Spectral Backbone]\label{inv:backbone}

\emph{Observation.}
A dominant subset of the spectrum carries the majority of the
multiplicity and admits a specific arithmetic decomposition.

\emph{Exact statement.}
The sum of rational parts of all distinct eigenvalues is
\begin{equation}\label{eq:backbone}
  \sum_{k=1}^{9} a_k = 88.
\end{equation}
This equals the sum of the nine distinct eigenvalues (by
Invariant~\ref{inv:phi}).  It decomposes as
\begin{equation}\label{eq:87}
  88 = 1 + 87 = 1 + 3 \times 29.
\end{equation}

There exists a dominant spectral backbone~$\mathcal{B}$ (the lower
six sectors, containing $91$~modes) such that
$\sum_{k \in \mathcal{B}} d_k \gg
 \sum_{k \notin \mathcal{B}} d_k$.

\emph{Minimal statement.}
$\sum a_k = 88 = 1 + 3 \times 29$.

\emph{Interpretation.}
The spectrum has a core subspace following the $S^3$ harmonic pattern
and a residual upper structure.  Nonlinear dynamics are expected to be
governed primarily by backbone modes.

\emph{Status:} Exact (arithmetic). \quad
\textbf{Type:} Algebraic (arithmetic decomposition).

\emph{Implication.}
The spectral backbone determines the dominant dynamical channels.
The backbone is defined purely spectrally and does not rely on any
geometric embedding.
\end{invariant}

% ── Invariant 7 ──────────────────────────────────────────
\begin{invariant}[Departure Structure]\label{inv:departure}

\emph{Observation.}
The residual eigenvalues outside the backbone form a structured set.

\emph{Exact statement.}
The spectrum decomposes as
$\mathcal{S} = \mathcal{B} \cup \mathcal{D}$, where
$\mathcal{B}$ is the backbone (lower six sectors, 91~modes) and
$\mathcal{D}$ is the departure (upper three sectors):
\begin{equation}\label{eq:29}
  |\mathcal{D}| = d_7 + d_8 + d_9 = 9 + 16 + 4 = 29.
\end{equation}
The departure is itself structured: $29 = 3^2 + 4^2 + 2^2$.

\emph{Verification.}
$91 + 29 = 120 = n$. \ $\checkmark$

\emph{Minimal statement.}
$d_7 + d_8 + d_9 = 29 = 3^2 + 4^2 + 2^2$; \; $91 + 29 = 120$.

\emph{Interpretation.}
The departure modes represent the deviation from the idealised
$S^3$ harmonic structure.  They correspond to finite-graph
corrections: the modes where the discrete 600-cell geometry
departs from a smooth 3-sphere.  Their structured decomposition
into perfect squares is consistent with an underlying
representation-theoretic origin.

\emph{Status:} Exact (count); interpretation is partially derived. \quad
\textbf{Type:} Combinatorial (structured residual).

\emph{Implication.}
The spectrum separates into idealised harmonic structure plus
structured correction terms.
\end{invariant}

\bigskip

The seven invariants described above do not constitute independent
assumptions, but rather appear as interrelated structural properties
of the $\Hfour$ Laplacian spectrum.  Their simultaneous presence
suggests that the spectrum is governed by a highly constrained
algebraic organisation rather than arising generically.


% ══════════════════════════════════════════════════════════
\section{Relation to Harmonic Structure on \texorpdfstring{$S^3$}{S³}}
\label{sec:harmonics}
% ══════════════════════════════════════════════════════════

The vertices of the 600-cell lie on the unit 3-sphere
$S^3 \subset \R^4$.  The Laplace--Beltrami operator on the continuous
$S^3$ has eigenvalues $\ell(\ell+2)$ with degeneracies $(\ell+1)^2$ for
$\ell = 0, 1, 2, \ldots$~\cite{vilenkin1993}.

The first six multiplicities of the 600-cell Laplacian ---
$1, 4, 9, 16, 25, 36$ --- exactly match $(\ell+1)^2$ for
$\ell = 0, 1, 2, 3, 4, 5$.  This is consistent with the interpretation
that the low-lying eigenmodes of the 600-cell graph correspond to
restrictions of $S^3$ hyperspherical harmonics to the 120-vertex
discretisation.

For $\ell \geq 6$ the pattern breaks.  A smooth $S^3$ would yield
degeneracies $49, 64, 81, \ldots$, but the finite graph instead
produces $9, 16, 4$.  The discrete structure cannot support harmonics
of arbitrarily high angular momentum; the upper modes fold back to
lower multiplicities (Invariant~\ref{inv:fold}).

We emphasise that this correspondence is structural rather than exact.
The 600-cell graph does not reproduce the full Laplace--Beltrami operator
on~$S^3$, but appears to encode a finite, symmetry-constrained
truncation.

This suggests that the 600-cell graph acts as a finite,
symmetry-reduced harmonic truncation of~$S^3$.  The truncation is not
arbitrary: it is governed by the $\Hfour$ symmetry and produces the
specific fold-back pattern described in Invariant~\ref{inv:fold}.

\begin{figure}[H]
\centering
\includegraphics[width=0.75\textwidth]{fig_harmonic_hint_s3.png}
\caption{Multiplicities of the 600-cell Laplacian compared with the
$(\ell+1)^2$ degeneracy pattern of $S^3$ hyperspherical harmonics.
The first six sectors match exactly; the upper three fold back.}
\label{fig:harmonics}
\end{figure}


% ══════════════════════════════════════════════════════════
\section{Algebraic Interpretation of the Invariants}
\label{sec:algebra}
% ══════════════════════════════════════════════════════════

The seven invariants are not independent observations; they arise from
a common algebraic source.

\subsection{Number-Theoretic Structure}

All eigenvalues lie in the number field $\Q(\sqrt{5}\,)$, the minimal
extension of the rationals containing the golden ratio.  The irrational
eigenvalues satisfy the minimal polynomials
\begin{align}
  x^2 - 18x + 36 &= 0 \qquad \text{(roots: } 9 \pm 3\sqrt{5}\text{)},
    \label{eq:minpoly1} \\
  x^2 - 20x + 80 &= 0 \qquad \text{(roots: } 10 \pm 2\sqrt{5}\text{)}.
    \label{eq:minpoly2}
\end{align}
The Galois group $\mathrm{Gal}(\Q(\sqrt{5}\,)/\Q) \cong \Z/2\Z$
acts by $\sqrt{5} \mapsto -\sqrt{5}$, which exchanges the members of
each conjugate pair while preserving eigenspace dimensions
(Invariant~\ref{inv:conj}).  The $\vphi$-cancellation
(Invariant~\ref{inv:phi}) is a necessary consequence of this Galois
symmetry applied to the trace.  Thus the spectrum is invariant under
the action of $\mathrm{Gal}(\Q(\sqrt{5}\,)/\Q)$.

\subsection{Graph-Theoretic Structure}

The 600-cell graph is vertex-transitive and has palindromic distance
shells $(1, 12, 32, 42, 32, 1)$ with diameter~5.  Numerical
investigation shows that the graph is \emph{not} distance-regular:
the intersection numbers are not constant at distances 2, 3, and~4.
Consequently the distance matrices do not close under multiplication,
and the graph does not admit a classical Bose--Mesner algebra of
dimension $d+1 = 6$.

This is consistent with the spectrum containing \emph{nine} distinct
eigenvalues rather than six (which would be expected for a
distance-regular graph of diameter~5).  The additional three eigenvalues
reflect the richer symmetry structure of the $\Hfour$ automorphism
group acting on the graph.  However, the palindromic shell structure
suggests a partial distance symmetry that is not captured by classical
distance-regular graph theory.

\subsection{Representation-Theoretic Perspective}

The adjacency matrix of the 600-cell commutes with the permutation
representation of~$\Hfour$ on $\R^{120}$.  The eigenspaces of the
Laplacian therefore decompose into irreducible representations
of~$\Hfour$.  The existence of nine distinct eigenvalues suggests that the
permutation representation of~$\Hfour$ on~$\R^{120}$ decomposes into
nine invariant subspaces, corresponding to nine irreducible components.
This decomposition is constrained by the commutation of the Laplacian
with the group action.

The multiplicities $d_k$ are the dimensions of these irreducible
representations.  The observation that the first six multiplicities
follow $(\ell+1)^2$ is consistent with the branching rules from the
continuous group $\mathrm{SO}(4)$ to the discrete subgroup~$\Hfour$.
A complete derivation of the irreducible decomposition from the
character table of~$\Hfour$ (which has 34~conjugacy classes) remains
an open problem.

The observed invariants suggest that the $\Hfour$ Laplacian spectrum
is governed by a highly constrained finite spectral algebra whose
structure is determined by the interplay between the non-crystallographic
Coxeter symmetry and the $\vphi$-geometry of the underlying polytope.


% ══════════════════════════════════════════════════════════
\section{Relation to the Dynamics of Part~I}\label{sec:dynamics}
% ══════════════════════════════════════════════════════════

The structural invariants identified in this paper are properties of
the static Laplacian spectrum, not dynamical claims.  However, they
provide a structural mechanism consistent with the dynamical phenomena
reported in Part~I~\cite{paper1}.

\begin{itemize}[nosep]
\item \textbf{Spectral confinement} follows from the small number of
  distinct eigenvalues (nine sectors rather than 120), which is itself
  a consequence of the $\Hfour$ symmetry
  (Invariants~\ref{inv:fold},~\ref{inv:zones}).

\item \textbf{Shell reflection} is enabled by the palindromic distance
  shells, which provide a symmetric propagation channel.  The spectral
  backbone (Invariant~\ref{inv:backbone}) determines the available
  resonance channels.

\item \textbf{Recurrent localisation} arises from the finite
  beat-frequency algebra: nine sectors produce only $\binom{9}{2} = 36$
  beat frequencies, which are sparse enough for periodic constructive
  interference.

\item \textbf{Breather formation} depends on the degenerate eigenspace
  structure (Invariant~\ref{inv:conj}), not merely on the eigenvalue
  spectrum, as demonstrated by the isospectral surrogate experiments
  of Part~I.
\end{itemize}

These invariants therefore provide a structural mechanism consistent
with the unusually coherent nonlinear dynamics observed on the $\Hfour$
shell.  The invariants constrain the space of possible dynamics, rather
than determining a unique dynamical outcome.  No dynamical assumptions
are required for these spectral constraints to hold.


% ══════════════════════════════════════════════════════════
\section{Discussion}\label{sec:discussion}
% ══════════════════════════════════════════════════════════

The structure identified here appears to be highly non-generic within
the space of 12-regular graphs on 120~vertices.  The Laplacian spectrum of the
600-cell graph is too structured to be generic.  A random 12-regular graph on 120~vertices would typically
produce 120~distinct eigenvalues with no algebraic relationships among
them.  The 600-cell produces nine, organized by conjugate pairing,
anchored at the degree, obeying a perfect-square sequence for six
sectors, and exhibiting exact $\vphi$-cancellation.

These invariants appear to arise from the exceptional position of the
600-cell at the terminus of the non-crystallographic Coxeter family.
The $H_n$ symmetry hierarchy terminates at $\Hfour$; there is no
analogous regular polytope in dimension five or higher.  The 600-cell
is therefore the most complex regular structure this symmetry family
can produce, and its Laplacian spectrum reflects this maximality.

Several questions remain open:
\begin{enumerate}[nosep]
\item Can the nine eigenvalues and their multiplicities be derived
  analytically from the character table of~$\Hfour$?
\item Does the fold-back pattern (Invariant~\ref{inv:fold}) have a
  representation-theoretic explanation in terms of branching rules from
  $\mathrm{SO}(4)$ to~$\Hfour$?
\item Is there a closed-form expression for the intersection array or
  an extended Bose--Mesner-like algebra that accounts for all nine
  eigenvalues?
\item Do the invariants extend or modify when the $\Hfour$ root system
  is embedded in the $\Eeight$ lattice (which decomposes as
  $240 = 2 \times 120$ roots)?
\end{enumerate}


% ══════════════════════════════════════════════════════════
\section{Conclusion}\label{sec:conclusion}
% ══════════════════════════════════════════════════════════

The 600-cell graph Laplacian spectrum exhibits a set of nontrivial
algebraic invariants that organize its nine eigenvalues and their
multiplicities into a highly constrained structure.  The
$\vphi$-cancellation, degree anchoring, conjugate pairing,
square-multiplicity folding, three-zone partition, spectral backbone,
and departure structure collectively demonstrate that this spectrum is
far more ordered than that of a generic graph with the same size and
degree.

These invariants arise from the exceptional symmetry of the $\Hfour$
Coxeter group and the $\vphi$-geometry of the 600-cell polytope.  They
motivate deeper work on the exact algebraic derivation of the spectrum
from the representation theory of~$\Hfour$, and they provide the
structural description consistent with the dynamical phenomena
reported in Part~I of this series.  The results motivate the search for
an exact algebraic derivation of the $\Hfour$ Laplacian spectrum from
representation-theoretic principles.  In particular, the coexistence of
Galois symmetry, square multiplicity structure, and spectral compression
in a single finite graph appears to be rare and warrants further
classification.

\paragraph{Reproducibility.}
All spectral computations and numerical verifications were performed
using explicit construction of the 600-cell graph and standard
eigendecomposition.  Code and data are available in the accompanying
repository.


% ══════════════════════════════════════════════════════════
% References
% ══════════════════════════════════════════════════════════
\begin{thebibliography}{10}

\bibitem{paper1}
L.~Smart,
``Symmetry-Constrained Resonant Modes on $H_4$ Geometry:
A Geometric Framework for Vibrational Field Dynamics,''
\textit{VFD Working Paper Series}, WO-VFD-PAPER-001, 2025.

\bibitem{humphreys1990}
J.~E.~Humphreys,
\textit{Reflection Groups and Coxeter Groups}.
Cambridge University Press, 1990.

\bibitem{coxeter1973}
H.~S.~M.~Coxeter,
\textit{Regular Polytopes}, 3rd~ed.
Dover Publications, 1973.

\bibitem{conway1999}
J.~H.~Conway and N.~J.~A.~Sloane,
\textit{Sphere Packings, Lattices and Groups}, 3rd~ed.
Springer, 1999.

\bibitem{vilenkin1993}
N.~J.~Vilenkin and A.~U.~Klimyk,
\textit{Representation of Lie Groups and Special Functions}, Vol.~2.
Kluwer Academic, 1993.

\bibitem{chung1997}
F.~R.~K.~Chung,
\textit{Spectral Graph Theory}.
AMS, 1997.

\bibitem{brouwer2012}
A.~E.~Brouwer and W.~H.~Haemers,
\textit{Spectra of Graphs}.
Springer, 2012.

\bibitem{serre1977}
J.-P.~Serre,
\textit{Linear Representations of Finite Groups}.
Springer, 1977.

\end{thebibliography}

\vfill
\noindent\rule{\textwidth}{0.4pt}\\[4pt]
{\footnotesize
\textit{Vibrational Field Dynamics Working Paper Series}
\hfill Part~II \quad Draft v1.1 (submission-ready)}

\end{document}
